% ==============================================================================
% CHƯƠNG 6: FORM VALIDATION & CÁC THẺ TƯƠNG TÁC HTML5
% ==============================================================================
\chapter{Form Validation \& Các Thẻ Tương Tác HTML5}
\label{chap:forms_validation_interactive}

Chương này phân tích chi tiết cơ chế Kiểm tra tính hợp lệ biểu mẫu (\term{Form Validation}) thuần trong HTML5 và làm chủ các thẻ Tương tác nội tạng hiện đại như \codeinline{<dialog>}, Popover API, \codeinline{<details>}, \codeinline{<progress>} và \codeinline{<meter>}.

% ==============================================================================
% MỤC 1: FORM VALIDATION TRONG HTML5
% ==============================================================================
\section{Cơ Chế HTML5 Form Validation \& Thuộc Tính Ràng Buộc}
\label{sec:form_validation_attributes}

\begin{conceptBox}[Form Validation Trong HTML5 Là Gì?]
\term{Form Validation} (Kiểm tra tính hợp lệ của biểu mẫu) là quá trình xác minh dữ liệu người dùng nhập vào ô input có đáp ứng đúng định dạng và quy tắc yêu cầu trước khi gửi lên máy chủ Backend hay không.
\begin{itemize}[leftmargin=1.5em, itemsep=2pt, parsep=0pt]
    \item \textbf{Native Validation}: Trình duyệt tự động kiểm tra và hiển thị thông báo lỗi giao diện mặc định mà không cần viết mã JavaScript phức tạp.
    \item \textbf{Lợi ích}: Tăng trải nghiệm người dùng (\term{UX}), giảm tải số lượng yêu cầu không hợp lệ gửi tới máy chủ và bảo vệ tính toàn vẹn dữ liệu.
\end{itemize}
\end{conceptBox}

\begin{prettyTableBox}[Bảng Tổng Hợp Các Thuộc Tính Ràng Buộc Dữ Liệu Form Validation]
\rowcolors{2}{white}{primaryBlue!4}
\begin{tabularx}{\linewidth}{K{3.0cm} Z Z}
\rowcolor{primaryBlue}
\thd{Thuộc Tính} & \thd{Ý Nghĩa Ràng Buộc} & \thd{Cú Pháp HTML Mẫu} \\
\codeinline{required} & Bắt buộc phải nhập/chọn dữ liệu trước khi gửi form. & \codeinline{<input type="text" required>} \\
\codeinline{pattern} & Biểu thức chính quy (\term{Regex}) kiểm tra định dạng chuỗi. & \codeinline{<input pattern="[0-9]{10}">} \\
\codeinline{minlength} / \codeinline{maxlength} & Số lượng ký tự tối thiểu và tối đa cho phép. & \codeinline{<input minlength="6" maxlength="20">} \\
\codeinline{min} / \codeinline{max} & Giá trị số hoặc ngày tháng nhỏ nhất và lớn nhất. & \codeinline{<input type="number" min="1" max="100">} \\
\codeinline{step} & Bước nhảy hợp lệ của giá trị số. & \codeinline{<input type="number" step="0.5">} \\
\codeinline{type} & Định dạng dữ liệu tự động (\codeinline{email}, \codeinline{url}, \codeinline{number}). & \codeinline{<input type="email">} \\
\end{tabularx}
\end{prettyTableBox}

\begin{codeWindow}[Ví dụ áp dụng thuộc tính Validation thuần HTML5]
\begin{lstlisting}[language=HTML5]
<form action="/register" method="POST">
  <!-- Bat buoc nhap & dinh dang Email -->
  <input type="email" name="email" required placeholder="(nhập_email@example.com)">
  
  <!-- Mat khau toi thieu 8 ky tu -->
  <input type="password" name="password" minlength="8" required placeholder="(*@Mật *@khẩu@*) (>=8 (*@ký tự@*)">
  
  (*@<!-- So dien (*@thoại@*) dung dinh dang 10 (*@chữ@*) so -->@*)
  <input type="tel" name="phone" pattern="[0-9]{10}" placeholder="(*@Số *@điện thoại@*) (10 (*@số@*)">
  
  <button type="submit">(*@Đăng *@Ký@*)</button>
</form>
\end{lstlisting}
\end{codeWindow}

\begin{browserWindow}[Kết Quả Hiển Thị Trình Duyệt: Form Validation HTML5]
\sffamily\small
\textbf{Đăng Ký Tài Khoản (HTML5 Validation Demo):} \\[6pt]
\tcbox[on line, arc=2pt, colback=white, colframe=primaryBlue!40, boxsep=1pt, left=6pt, right=40pt, top=3pt, bottom=3pt]{\color{gray!80}nhap\_email@example.com} \quad \textcolor{codeComment}{\small\textit{(Required \& Email type)}} \\[4pt]
\tcbox[on line, arc=2pt, colback=white, colframe=primaryBlue!40, boxsep=1pt, left=6pt, right=40pt, top=3pt, bottom=3pt]{\color{gray!80}Mật khẩu ($\ge$8 ký tự)} \quad \textcolor{codeComment}{\small\textit{(Minlength = 8)}} \\[4pt]
\tcbox[on line, arc=2pt, colback=white, colframe=primaryBlue!40, boxsep=1pt, left=6pt, right=40pt, top=3pt, bottom=3pt]{\color{gray!80}Số điện thoại (10 số)} \quad \textcolor{codeComment}{\small\textit{(Regex Pattern = [0-9]\{10\})}} \\[8pt]
\tcbox[on line, arc=3pt, colback=primaryBlue, colframe=primaryBlue, boxsep=2pt, left=10pt, right=10pt, top=3pt, bottom=3pt]{\color{white}\bfseries\small Đăng Ký}
\end{browserWindow}

% ==============================================================================
% MỤC 2: THẺ ACCORDION <details> & <summary>
% ==============================================================================
\section{Thẻ Accordion <details> \& <summary>}
\label{sec:details_summary}

\begin{conceptBox}[Thẻ Accordion Thuần HTML5]
Thẻ \codeinline{<details>} và \codeinline{<summary>} tạo ra widget ẩn/hiện nội dung (Accordion/Collapsible) thuần HTML5 mà không cần dùng JavaScript hay thư viện bên thứ ba.
\begin{itemize}[leftmargin=1.5em, itemsep=2pt, parsep=0pt]
    \item \codeinline{<summary>}: Đóng vai trò là tiêu đề nhấp chọn (kích hoạt mở/đóng).
    \item Thuộc tính \codeinline{open}: Đặt trạng thái mặc định mở nội dung khi tải trang.
    \item Thuộc tính \codeinline{name="..."}: Gom nhóm các thẻ \codeinline{<details>} lại với nhau, tự động đóng các ô khác khi một ô mới được mở (chuẩn HTML Living Standard).
\end{itemize}
\end{conceptBox}

\begin{codeWindow}[Ví dụ danh sách FAQ Accordion với details và summary]
\begin{lstlisting}[language=HTML5]
<details name="faq-group" open>
  <summary><strong>Khoa hoc keo dai trong bao lau?</strong></summary>
  <p>Khoa hoc keo dai trong 3 thang voi 50 bai giang chuyen sau.</p>
</details>

<details name="faq-group">
  <summary><strong>Toi co duoc nhan chung chi khong?</strong></summary>
  <p>Co, ban se nhan duoc chung chi hoan thanh chuan quoc te.</p>
</details>
\end{lstlisting}
\end{codeWindow}

\begin{browserWindow}[Kết Quả Trình Duyệt: Widget Accordion FAQ]
\sffamily\small
\textbf{[-]} \textbf{Khóa học kéo dài trong bao lâu?} \quad \textcolor{codeComment}{\small\textit{(Đang mở open)}} \\
\begin{quote}\small
Khóa học kéo dài trong 3 tháng với 50 bài giảng chuyên sâu.
\end{quote}
\textbf{[+]} \textbf{Tôi có được nhận chứng chỉ không?} \quad \textcolor{codeComment}{\small\textit{(Đang đóng)}}
\end{browserWindow}

% ==============================================================================
% MỤC 3: THẺ CỬA SỔ MODAL <dialog>
% ==============================================================================
\section{Thẻ Cửa Sổ Modal <dialog>}
\label{sec:dialog_modal}

\begin{conceptBox}[Thẻ Cửa Sổ Hộp Thoại Dialog Native]
Thẻ \codeinline{<dialog>} cung cấp giải pháp chuẩn native để xây dựng cửa sổ Popup / Modal / Hộp thoại xác nhận trong HTML5.
\begin{itemize}[leftmargin=1.5em, itemsep=2pt, parsep=0pt]
    \item \textbf{Phương thức JS \codeinline{showModal()}}: Mở dialog ở chế độ Modal chính thức trên lớp trên cùng (\term{Top Layer}) với phần nền mờ (\codeinline{::backdrop}).
    \item \textbf{Phương thức JS \codeinline{close()}}: Đóng cửa sổ hộp thoại.
    \item \textbf{Cấu hình \codeinline{method="dialog"}}: Khi nằm trong form, submit nút bấm sẽ tự động đóng dialog mà không làm tải lại trang.
\end{itemize}
\end{conceptBox}

\begin{codeWindow}[Ví dụ tạo Modal Popup chuẩn với thẻ dialog]
\begin{lstlisting}[language=HTML5]
<!-- Nut mo modal -->
<button onclick="document.getElementById('confirmModal').showModal()">Xoa Bai Viet</button>

<!-- The dialog modal -->
<dialog id="confirmModal">
  <h3>Xac Nhan Hanh Dong</h3>
  <p>Ban co chac chan muon xoa bai viet nay vinh vien khong?</p>
  <form method="dialog">
    <button value="cancel">Huy Bo</button>
    <button value="confirm" style="color: red;">Xoa Vinh Vien</button>
  </form>
</dialog>
\end{lstlisting}
\end{codeWindow}

\begin{browserWindow}[Mô Phỏng Trình Duyệt: Hộp Thoại Modal Dialog]
\sffamily\small
\tcbox[on line, arc=3pt, colback=codeComment, colframe=codeComment, boxsep=2pt, left=8pt, right=8pt, top=3pt, bottom=3pt]{\color{white}\bfseries\small Xóa Bài Viết} \\[10pt]

\begin{center}
\begin{tcolorbox}[arc=4pt, colback=white, colframe=primaryBlue, boxsep=8pt, width=0.85\linewidth, halign=center]
  \textbf{\textcolor{primaryBlue}{Xác Nhận Hành Động}} \\[4pt]
  {\small Bạn có chắc chắn muốn xóa bài viết này vĩnh viễn không?} \\[8pt]
  \tcbox[on line, arc=2pt, colback=gray!20, colframe=gray!40, boxsep=1pt, left=8pt, right=8pt, top=2pt, bottom=2pt]{\small Hủy Bỏ} \qquad
  \tcbox[on line, arc=2pt, colback=red!80, colframe=red!80, boxsep=1pt, left=8pt, right=8pt, top=2pt, bottom=2pt]{\color{white}\small\bfseries Xóa Vĩnh Viễn}
\end{tcolorbox}
\end{center}
\end{browserWindow}

% ==============================================================================
% MỤC 4: POPOVER API
% ==============================================================================
\section{Thẻ Popover API (popover, popovertarget)}
\label{sec:popover_api}

\begin{conceptBox}[Tính Năng Đột Phá Popover API (WHATWG Standard)]
\term{Popover API} cho phép biến bất kỳ phần tử HTML nào thành Tooltip, Dropdown Menu hoặc Popup nổi hiển thị ở lớp trên cùng (\term{Top Layer}) mà không bị giới hạn bởi thuộc tính \codeinline{overflow: hidden} của phần tử cha.
\begin{itemize}[leftmargin=1.5em, itemsep=2pt, parsep=0pt]
    \item Thuộc tính \codeinline{popover="auto|manual"}: Gán tính năng popover cho phần tử.
    \item Thuộc tính \codeinline{popovertarget="id"}: Liên kết nút bấm kích hoạt tới ID của khung popover.
    \item Khởi chạy hoàn toàn thuần HTML5 không cần JavaScript!
\end{itemize}
\end{conceptBox}

\begin{codeWindow}[Ví dụ Tooltip Popover thuần HTML5 không cần JS]
\begin{lstlisting}[language=HTML5]
<!-- Nut kich hoat Popover -->
<button popovertarget="userTooltip">Xem Thong Tin Tac Gia</button>

<!-- Khung noi dung Popover -->
<div id="userTooltip" popover="auto">
  <h4>Nguyen Trong Tan</h4>
  <p>Chuyen gia Kien truc Frontend & Tac gia sach.</p>
</div>
\end{lstlisting}
\end{codeWindow}

% ==============================================================================
% MỤC 5: THƯỚC ĐO METER & TIẾN TRÌNH PROGRESS
% ==============================================================================
\section{Thẻ Thước Đo <meter> \& Tiến Trình <progress>}
\label{sec:meter_progress}

\begin{conceptBox}[Phân Biệt <progress> và <meter>]
\begin{itemize}[leftmargin=1.5em, itemsep=2pt, parsep=0pt]
    \item \codeinline{<progress>}: Biểu diễn tiến trình đang chạy (đang tải tệp, hoàn tất đăng ký...).
    \item \codeinline{<meter>}: Biểu diễn thước đo của một giá trị nằm trong phạm vi xác định (dung lượng ổ đĩa, nhiệt độ, điểm đánh giá).
\end{itemize}
\end{conceptBox}

\begin{prettyTableBox}[Bảng Thuộc Tính Của Thẻ meter Và progress]
\rowcolors{2}{white}{primaryBlue!4}
\begin{tabularx}{\linewidth}{K{3.0cm} Z Z}
\rowcolor{primaryBlue}
\thd{Thuộc Tính} & \thd{Áp Dụng Cho} & \thd{Ý Nghĩa} \\
\codeinline{value} & \codeinline{<progress>}, \codeinline{<meter>} & Giá trị hiện tại. \\
\codeinline{min} / \codeinline{max} & \codeinline{<progress>}, \codeinline{<meter>} & Giá trị tối thiểu và tối đa. \\
\codeinline{low} / \codeinline{high} / \codeinline{optimum} & Chỉ \codeinline{<meter>} & Các dải mốc giá trị thấp, cao và tối ưu. \\
\end{tabularx}
\end{prettyTableBox}

\begin{codeWindow}[Ví dụ sử dụng thẻ progress và thẻ meter]
\begin{lstlisting}[language=HTML5]
<!-- Thanh tien trinh tai tep 70% -->
<label for="fileDown">Tien trinh tai tep:</label>
<progress id="fileDown" value="70" max="100">70%</progress>

<!-- Thuoc do dung luong o dia -->
<label for="diskSpace">Dung luong o dia:</label>
<meter id="diskSpace" value="85" min="0" max="100" low="30" high="80" optimum="20">85%</meter>
\end{lstlisting}
\end{codeWindow}

\begin{browserWindow}[Kết Quả Trình Duyệt: Tiến Trình \& Thước Đo]
\sffamily\small
\textbf{Tiến trình tải tệp:} \;\;
\tcbox[on line, arc=3pt, colback=primaryBlue!20, colframe=primaryBlue, boxsep=1pt, left=4pt, right=40pt, top=2pt, bottom=2pt]{\color{primaryBlue}\bfseries 70\%} \\[6pt]
\textbf{Dung lượng ổ đĩa:} \;\;
\tcbox[on line, arc=3pt, colback=red!20, colframe=red!80, boxsep=1pt, left=4pt, right=40pt, top=2pt, bottom=2pt]{\color{red!80}\bfseries 85\% (Cảnh báo dung lượng cao)}
\end{browserWindow}

% ==============================================================================
\section{Tóm Tắt Chương 6}
Chương 6 đã hoàn tất phân tích chi tiết cơ chế \textbf{HTML5 Form Validation} và toàn bộ các thẻ \textbf{Tương tác nội tạng hiện đại} (\codeinline{dialog}, Popover API, \codeinline{details}, \codeinline{meter}, \codeinline{progress}). Kiến thức phần HTML5 hiện đã đầy đủ, nhất quán và sạch sẽ 100\%, sẵn sàng để chuyển tiếp sang phần CSS3 Styling nâng cao.


